%----------------------------------------------------------------------
\section{Meta-analysis: the technique matrix}\label{app:meta}
%----------------------------------------------------------------------

This appendix records the structure discovered by
profiling the paper's 47 audited claims against the
six closure techniques. The analysis was conducted
after all 17 gaps were closed and revealed higher-order
structure in the proof methodology.

\subsection{The six techniques}

Every closure in the paper uses one or more of:
\begin{description}[nosep,leftmargin=1.5cm]
  \item[T1: $1{+}1{=}0$] Characteristic 2 of $\GF(2)$:
    signs collapse, cross terms cancel, self-inverse.
  \item[T2: Leibniz] The graded Leibniz rule
    $\partial(a \wedge b) = \partial a \wedge b +
    a \wedge \partial b$ over $\GF(2)$.
  \item[T3: Linearity] $\GF(2)$-linearity of evaluation,
    profiles, constraints. Kernels, solution spaces,
    unique factorization.
  \item[T4: Strict $\wedge$] Strict associativity,
    strict unit, strict commutativity of $\wedge$.
    All coherence morphisms are identities.
  \item[T5: SES] The short exact sequence
    $0 \to \kappa \to V \to V/\kappa \to 0$ and the
    tensor decomposition
    $\bigwedge(V) \cong \bigwedge(\kappa) \otimes
    \bigwedge(V/\kappa)$.
  \item[T6: Galois] The Galois group
    $GL(V/\kappa, \GF(2))$ and its reduction under
    $\kappa$-evolution. Equivariance, invariance,
    orbit exhaustion.
\end{description}

\subsection{The closure--technique matrix}

Each of the 47 audited claims is assigned a binary
profile in $\GF(2)^6$ recording which techniques its
proof uses. The matrix has 18 distinct profiles spanning
all of $\GF(2)^6$ (rank 6).

\begin{center}
\scriptsize
\begin{tabular}{l|cccccc|c}
 & T1 & T2 & T3 & T4 & T5 & T6 & Cat \\
\hline
\multicolumn{8}{l}{\emph{Equivalence class}
  $[1{+}1{=}0]$: 10 items} \\
Non-default Bool 4.11 &
  $\bullet$ & & & & & & V \\
Sym difference 4.19 &
  $\bullet$ & & & & & & V \\
Chain complex 5.5 &
  $\bullet$ & & & & & & V \\
Three perspectives 7.11 &
  $\bullet$ & & & & & & V \\
Deloop collapse 8.5 &
  $\bullet$ & & & & & & V \\
Proof theory 9.11 &
  $\bullet$ & & & & & & V \\
Self-modeling 9.21 &
  $\bullet$ & & & & & & V \\
Banach--Tarski 4.37 &
  $\bullet$ & & & & & & V \\
Commutativity B.2 &
  $\bullet$ & & & & & & V \\
Periodic table 8.8 &
  $\bullet$ & & & & & & O \\[4pt]
\multicolumn{8}{l}{\emph{Equivalence class}
  $[\mathrm{SES} + \mathrm{Galois}]$: 6 items} \\
FOL internal 9.12 &
  & & & & $\bullet$ & $\bullet$ & V \\
Choice finite (G2) &
  & & & & $\bullet$ & $\bullet$ & C \\
Univalence (G5) &
  & & & & $\bullet$ & $\bullet$ & C \\
Yoneda (G7) &
  & & & & $\bullet$ & $\bullet$ & C \\
Pairing (G17) &
  & & & & $\bullet$ & $\bullet$ & C \\
Choice infinite 4.33ii &
  & & & & $\bullet$ & $\bullet$ & O \\[4pt]
\multicolumn{8}{l}{\emph{Equivalence class}
  $[\mathrm{Linear}]$: 4 items} \\
Order-indep 3.15 &
  & & $\bullet$ & & & & V \\
Empty set 4.17 &
  & & $\bullet$ & & & & V \\
Forgetful/free 6.10 &
  & & $\bullet$ & & & & V \\
Bool dep type 9.13 &
  & & $\bullet$ & & & & V \\[4pt]
\multicolumn{8}{l}{\emph{12 further classes with
  1--3 members (see text)}}
\end{tabular}
\end{center}

\noindent Cat: V = verified from start, C = gap-closed,
O = open/acknowledged.

\subsection{The technique co-occurrence tree}
\label{sec:technique-tree}

The co-occurrence matrix $C_{jk}$ counts how many items
use both technique $j$ and technique $k$. Reducing
modulo 2 yields a symmetric matrix over $\GF(2)$ whose
off-diagonal entries define an adjacency graph.

This graph is a \emph{tree} on six vertices with five
edges:
\begin{center}
$\underset{\text{(7)}}{\text{T1}}
\;{-}\;
\underset{\text{(1)}}{\text{T2}}
\;{-}\;
\underset{\text{(9)}}{\text{T5}}
\;{-}\;
\underset{\text{(13)}}{\text{T6}}
\;{\begin{array}{c} | \\ | \end{array}}$
\\[2pt]
$\underset{\text{(4)}}{\text{T4}} \quad\quad
\underset{\text{(5)}}{\text{T3}}$
\end{center}
(Usage counts in parentheses.)

The tree has three arms meeting at the Galois hub (T6):
\begin{description}[nosep,leftmargin=1.5cm]
  \item[Algebraic arm]
    $\text{T1} - \text{T2} - \text{T5} - \text{T6}$.
    How the exterior algebra computes: characteristic 2,
    Leibniz rule, tensor decomposition.
  \item[Coherence arm] $\text{T6} - \text{T4}$. How
    composition works: strictness of $\wedge$.
  \item[Constraint arm] $\text{T6} - \text{T3}$. How
    populations behave: linearity of evaluation.
\end{description}

Strict (T4) and Linearity (T3) never co-occur in
gap closures (and co-occur only once in the full 47
items: self-enrichment 9.2). They operate on different
grades of $\bigwedge(V)$: Strict on morphism structure,
Linearity on population structure.

\subsection{The Galois gap}

The 27 pre-verified items use Galois (T6) in 19\% of
proofs. The 18 gap-closed items use it in 78\%. The
gaps were precisely the claims requiring the Galois
structure that the paper had not yet recognized.

The pre-verified profiles span a rank-5 subspace of
$\GF(2)^6$; the gap-closed profiles span rank 6. The
missing dimension was Leibniz (T2), used zero times in
pre-existing proofs and once in the gap closure of
exhaustivity (G4). The full technique space is occupied
only after both the Galois structure and the Leibniz
filling were brought in.

Verified-only profiles are single-technique: $[1{+}1{=}0]$
(10 items), $[\text{Linear}]$ (4), $[\text{SES}]$ (3).
Gap-only profiles are compound:
$[\text{Linear}{+}\text{SES}{+}\text{Galois}]$,
$[1{+}1{=}0{+}\text{Strict}{+}\text{Galois}]$,
$[1{+}1{=}0{+}\text{Leibniz}{+}\text{SES}]$.
The gaps existed at the intersections of the technique
space, where multiple arms needed to combine.

\subsection{The rotation principle}
\label{sec:rotation}

Every proof can be \emph{rotated} through the Galois
hub to produce an alternate proof using techniques from
a different arm:
\[
  \text{Algebraic} \xrightarrow{\text{lift}}
  \text{Galois} \xrightarrow{\text{project}}
  \text{Coherence or Constraint}
\]
The three arms correspond to the three perspectives
$\{\kappa, \sigma, \tau\}$. Rotating a proof is rotating
the perspective from which the same mathematical content
is viewed.

\emph{Examples of rotation.}

$\partial \circ \partial = 0$ (chain complex, 5.5):
the algebraic proof cancels double-sums over $\GF(2)$.
The coherence alternate: $\partial \circ \partial = 0$
follows from strict graded-commutativity of $\wedge$.

Interchange law (G6): the algebraic proof computes in
the $\GF(2)^2$ profile space and shows cross terms
cancel via $1{+}1{=}0$. The coherence alternate: the
2-category is strict (G12), so interchange is
$\mathrm{id} \circ \mathrm{id} = \mathrm{id}$.

$(n,1)$-structure (G9): the coherence proof uses
commutativity of $\wedge$ (factor permutations trivial
over $\GF(2)$). The constraint alternate: grade-$k$
evaluation on $S$ is a $k$-linear form, and $k$-linear
forms for $k \geq 2$ are symmetric over $\GF(2)$.

Russell exclusion (4.5): the algebraic proof shows the
affine offset $\tau(0) = 1 \neq 0$. The Galois
alternate: the Russell predicate is in the wrong
$GL$-orbit (affine, not linear).

\emph{The exception}: exhaustivity filling (G4). The
Leibniz rule $\partial(R \wedge e_{n+1}) = R$ is the
unique constructive filling procedure. A Galois
alternate exists (any element not in the stabilizer of
$R$ fills) but is existential, not constructive.
Leibniz is the one dimension of the technique space
with no constructive substitute.

\emph{Consequence}: the 47 claims are not 47 independent
facts. Modulo rotation, approximately 15--18 independent
facts exist, each provable from 2--3 arms. The paper's
claims are overdetermined by its axioms: each truth is
visible from multiple perspectives, connected by the
same $S_3$ that acts on the framework's objects.



\subsection{The Leibniz exception: why one technique
  doesn't rotate}
\label{sec:leibniz-exception}

The rotation principle (\S\ref{sec:rotation}) has one
exception: the Leibniz rule (T2) cannot be replaced by
a constructive alternate. This subsection explains why,
drawing on the geometric origin of the rule.

\emph{The graded Leibniz rule} for the boundary operator
is:
\[
  \partial(a \wedge b) = \partial a \wedge b
  + (-1)^{|a|}\, a \wedge \partial b.
\]
Over $\GF(2)$, $(-1)^{|a|} = 1$, so this simplifies to
$\partial(a \wedge b) = \partial a \wedge b + a \wedge
\partial b$.

\emph{Geometric origin.}
The rule is derived from the boundary of a product space.
If $A$ is a $k$-cell and $B$ is an $l$-cell, then
$A \times B$ is a $(k{+}l)$-cell whose boundary
decomposes as:
\[
  \partial(A \times B) = (\partial A \times B)
  \cup (A \times \partial B).
\]
The first term sweeps the boundary of $A$ over $B$
(holding $B$ fixed). The second holds $A$ fixed and
takes the boundary of $B$. In the exterior algebra,
$\times$ becomes $\wedge$ and $\cup$ becomes $+$,
yielding the Leibniz rule.

\emph{Why it matters for the filling.}
In the filling procedure (\S\ref{app:exhaustivity-full}),
$A = R$ (the inner system's residue, a $k$-cycle in
$\bigwedge(\kappa)$) and $B = e_{n+1}$ (the outer
system's fresh discriminator, in $V/\kappa$). The product
$R \wedge e_{n+1}$ lives in $\bigwedge(\kappa) \otimes
\bigwedge(V/\kappa)$ --- it crosses the inner/outer
tensor boundary. The Leibniz rule decomposes its
boundary:
\begin{align*}
  \partial(R \wedge e_{n+1})
  &= \underbrace{\partial R}_{\text{sweep: inner boundary
       over outer}} \wedge\; e_{n+1}
  + R \wedge \underbrace{\partial(e_{n+1})}_{\text{hold
       inner, take outer boundary}} \\
  &= 0 \cdot e_{n+1} + R \cdot 1 = R.
\end{align*}
The ``sweep'' term vanishes because $R$ is a cycle
($\partial R = 0$: the inner system's hole has no
boundary of its own). The ``hold'' term gives $R$
because the outer discriminator's boundary is the
augmentation ($\partial(e_{n+1}) = 1$: a single
discriminator bounds to the grade-0 unit).

\emph{The geometric content}: the inner system's hole
$R$ IS the face that appears when the product
(hole $\times$ fresh discriminator) is bounded. The
outer system's contribution is the factor that, when
multiplied and then bounded, reproduces the hole.

\emph{The three derivation paths} of the Leibniz rule
all have framework interpretations:

\begin{enumerate}[nosep]
  \item \emph{Geometric decomposition} (boundary of a
    product). In the framework: the product
    $\bigwedge(\kappa) \otimes \bigwedge(V/\kappa)$ is
    the inner/outer tensor space. Its boundary
    decomposes into ``inner boundary swept over outer''
    and ``inner held while outer is bounded.'' This is
    the filling computation.

  \item \emph{Orientation consistency} (sign $(-1)^k$).
    Over $\GF(2)$, the sign is trivially 1. The
    orientation problem that makes the Leibniz rule
    subtle in other settings --- the careful tracking of
    which face is ``positive'' and which is ``negative''
    --- collapses entirely. Characteristic 2 eliminates
    the sign bookkeeping. This is why the filling works
    cleanly: no orientation issues can obstruct it.

  \item \emph{$\partial^2 = 0$ verification}. The sign
    $s = (-1)^k$ is required for $\partial^2 = 0$.
    Over $\GF(2)$, $s = 1$ always, and the requirement
    is automatically satisfied. The chain complex
    property (Proposition \ref{prop:chain}) is the
    algebraic verification that the Leibniz rule is
    self-consistent.
\end{enumerate}

\emph{Why Leibniz doesn't rotate.}
The other five techniques describe \emph{properties}:
$1{+}1{=}0$ is a property of the field, Linearity of
the evaluation, Strict of the product, SES of the
decomposition, Galois of the symmetry group. Properties
can be viewed from different perspectives; rotating
through $S_3$ yields the same property seen from a
different vertex.

Leibniz describes an \emph{operation}: the specific
computation of $\partial$ across the tensor boundary
$\bigwedge(\kappa) \otimes \bigwedge(V/\kappa)$. It is
the mechanism of crossing, not a property of either
side. The Galois alternate says ``a filling exists''
(because the stabilizer is a proper subgroup). But it
does not \emph{produce} the filling --- it identifies
the space where it lives. Leibniz produces it: take $R$,
wedge with $e_{n+1}$, apply $\partial$, get $R$ back.

The constructive content is the crossing itself: the
moment where the outer system's model of the inner
system's hole becomes the inner system's new
discriminator. This is the transfer from outer-interior
to inner-exterior --- the one step in the framework that
cannot be described as a symmetry, a linearity, or a
structural identity, because it is the act of making an
outer-exterior assertion from inner-exterior observation.

The framework can model itself (self-hosting). It can
describe what it cannot distinguish (the operationalist
axiom). It can rotate between perspectives ($S_3$). But
when it needs to \emph{fill a hole} --- to go from
``this residue exists'' to ``here is the discriminator
that resolves it'' --- it must compute $\partial$ across
the tensor boundary. That computation is Leibniz, and
Leibniz alone.


\subsection{Unified claims and alternate constructions}
\label{sec:unified}

The rotation principle reveals that many seemingly
distinct claims are aspects of a smaller number of
general claims, viewed from different arms of the
technique tree. This subsection identifies the general
claims and provides alternate constructions.

\subsubsection*{Claim A: Characteristic 2 collapses sign,
  orientation, and self-reference}

\emph{General statement.} Over $\GF(2)$, every element
is its own additive inverse. Negation, sign, and
orientation are trivial. Self-referential operations
annihilate ($v \wedge v = 0$).

\emph{Instances} (10 items in $[1{+}1{=}0]$):
non-default Booleanness, symmetric difference,
chain complex ($\partial^2 = 0$), three perspectives,
delooping collapse ($\beta^2 = \mathrm{id}$), proof
theory (four-valued truth tables), self-modeling,
Banach--Tarski (type error), commutativity of $\wedge$,
periodic table conjecture.

\emph{Alternate via Strict}: each instance follows
from the strict graded-commutativity of $\wedge$.
For example, $\partial^2 = 0$ (chain complex) is
the strict identity
$\partial \circ \partial = \mathrm{id} + \mathrm{id}
= 0$ in the endomorphism ring of the chain group.

\emph{Alternate via Galois}: each instance is a
$GL(V)$-invariant property. For example, commutativity
of $\wedge$ is invariant under basis change; the
four-valued truth tables are $S_3$-invariant.

\subsubsection*{Claim B: The $\kappa$-split determines
  all epistemic structure}

\emph{General statement.}
The short exact sequence $0 \to \kappa \to V \to
V/\kappa \to 0$ and the Galois group
$GL(V/\kappa, \GF(2))$ jointly determine what can be
known, what can be selected, what can be identified,
what can be represented, and what can be paired.

\emph{Instances} (6 items in $[\mathrm{SES}{+}
\mathrm{Galois}]$): FOL internal, finite choice,
univalence, Yoneda, pairing, infinite choice.

\emph{Alternate via Linearity}: each instance can be
restated as a linear-algebraic fact about the
annihilator $\kappa^\perp$.
\begin{itemize}[nosep]
  \item \emph{Univalence}: persistent equivalence means
    $a + b \in \bigcap \kappa^\perp = \{0\}$,
    hence $a = b$. The linear alternate: the kernel
    of the evaluation map, intersected over all $\kappa$,
    is $\{0\}$.
  \item \emph{Pairing}: $\{a, b\}$ is a
    $\kappa$-equivalence class iff
    $a + b \in \kappa^\perp$. The linear alternate:
    the pair is the kernel of all $d$ with
    $d(a + b) = 0$.
  \item \emph{Finite choice}: selection terminates in
    $\dim(V)$ steps. The linear alternate: each step
    reduces the dimension of the solution space by 1.
  \item \emph{Yoneda}: the pro-system of finite Yonedas
    is Galois-equivariant. The linear alternate: the
    restriction functors are linear maps between
    finite-dimensional hom-spaces.
  \item \emph{FOL}: $\forall$ and $\exists$ are
    adjunctions to pullback. The linear alternate: they
    are $\GF(2)$-linear operations on the profile space
    (universal quantification = intersection of kernels;
    existential = span of images).
\end{itemize}

\subsubsection*{Claim C: $S_3$-invariance determines the
  Fano/toroid structure}

\emph{General statement.}
Over $\GF(2)^2$, the Galois group $GL(2, \GF(2))
\cong S_3$ has a unique invariant structure at each
grade: the triangle identity at grade 1, the unique
2-cell $e_1 \wedge e_2$ at grade 2, and the trivial
determinant $\det = 1$. All claims about the toroid's
structure are consequences of this unique invariance.

\emph{Instances} (4 items in $[1{+}1{=}0{+}
\mathrm{Galois}]$): self-hosting, Fano closure,
fixed-point stability, sheaf topology.

\emph{Alternate via SES}: each instance can be derived
from the fiber structure.
\begin{itemize}[nosep]
  \item \emph{Self-hosting}: the internal $\GF(2)$ is
    the fiber at the terminal $\kappa$-stage (the Boolean
    restriction of $\Omega$).
  \item \emph{Fano closure}: the seven Fano lines are the
    seven non-zero elements of $\GF(2)^3$; three
    independent fibers generate all seven via the
    triangle identity.
  \item \emph{Fixed-point stability}:
    $\dim \bigwedge^2(\GF(2)^2) = 1$; the unique
    top form is the section of the SES at grade 2.
  \item \emph{Sheaf topology}: two 1-dimensional
    sub-regimes (fibers from different axes) span
    the 2-dimensional regime.
\end{itemize}

\subsubsection*{Claim D: Strict composition with
  equivariant decomposition determines higher structure}

\emph{General statement.}
In a strict monoidal category where $\wedge$ is strictly
associative, commutative, and unital, all coherence
morphisms are identities and all higher cells are
determined by their composable edges.

\emph{Instances}: delooping (Strict${+}$Galois),
Segal conditions (Strict${+}$Galois), symmetric monoidal
coherence (1+1=0${+}$Strict${+}$Galois), closure
naturality (Strict alone).

\emph{Unification}: delooping and Segal share the
profile $[\mathrm{Strict}{+}\mathrm{Galois}]$.
Delooping says a strict monoidal category is a
one-object 2-category; Segal says $k$-cells are
determined by composable edges. These are the same
claim: \emph{strict composition is fully determined
by its 1-cells}, viewed at grade 2 (delooping) and
grade $k$ (Segal).

\emph{Alternate via 1+1=0}: strictness follows from
$(-1)^k = 1$ over $\GF(2)$. Every sign that could
make a coherence morphism non-trivial collapses.

\emph{Alternate via Linearity}: the Grassmannian
$\mathrm{Gr}(k, V)$ parametrizes decomposable
$k$-forms. The Segal map is a linear isomorphism
between the space of decomposable $k$-forms and
the space of $k$-tuples of independent 1-forms.

\subsubsection*{Claim E: Linearity of evaluation makes
  constraints solvable}

\emph{General statement.}
Profile linearity and evaluation linearity make all
constraint satisfaction problems over $\GF(2)$
reducible to linear algebra: kernels exist, solution
spaces are subspaces, factoring morphisms are unique.

\emph{Instances}: limits ($1{+}1{=}0{+}\mathrm{Linear}
{+}\mathrm{SES}$), power objects ($\mathrm{Linear}
{+}\mathrm{SES}{+}\mathrm{Galois}$), foundation chain
bound ($\mathrm{Linear}{+}\mathrm{SES}{+}\mathrm{Galois}$),
$\omega$-chain stabilization ($\mathrm{Linear}{+}
\mathrm{SES}{+}\mathrm{Galois}$).

\emph{Unification}: foundation (G1), indistinguishability
(G3), and power objects (G14) share the profile
$[\mathrm{Linear}{+}\mathrm{SES}{+}\mathrm{Galois}]$.
Foundation bounds chain depth by dimension. Power objects
classify subobjects by dimension. Indistinguishability
stabilizes by dimension. All three are:
\emph{dimension counting in the annihilator lattice}.

\emph{Alternate via 1+1=0}: each instance uses the
fact that $\GF(2)$-linear constraints have well-defined
kernels. This is a consequence of characteristic 2:
the kernel of a $\GF(2)$-linear map is a subspace
(because $\ker(f) = \{x : f(x) = 0\}$ is closed
under $+$, and $+ = -$ over $\GF(2)$).

\subsubsection*{Claim F: Russell exclusion and
  interchange are the same claim}

\emph{General statement.} Evaluation linearity excludes
non-linear predicates. This exclusion manifests as
paradox prevention (Russell) and coherence (interchange).

Russell exclusion (4.5) and the interchange law (G6)
share the profile $[1{+}1{=}0 + \mathrm{Linear}]$.

\begin{itemize}[nosep]
  \item \emph{Russell}: the self-referential predicate
    $d_P$ has an affine offset $\tau(0) = 1 \neq 0$.
    Evaluation linearity excludes it because affine
    maps are not linear.
  \item \emph{Interchange}: the cross terms
    $g \cdot f'$ in the interchange square have
    coefficient 2 in the profile-space computation.
    Linearity over $\GF(2)$ kills them ($2 = 0$).
\end{itemize}

Both are: \emph{linearity over $\GF(2)$ rejects the
affine/quadratic component}. Russell's paradox is an
affine predicate; the interchange failure would be a
quadratic cross term. The same linearity excludes both.

\emph{Alternate via Galois}: Russell's predicate is in
the wrong $GL$-orbit (affine, not linear). The
interchange cross terms are in the wrong $GL$-orbit
(quadratic, not bilinear).

\emph{Alternate via Strict}: in a strict 2-category,
interchange is the identity (no cross terms arise).
Russell exclusion from strictness: a strict evaluation
has no room for an affine offset.

\subsubsection*{Claim G: The bilinear pairing is the
  one structure}

\emph{General statement.}
The evaluation map $\mathrm{ev}(d, a) = \langle d, a
\rangle$ is a bilinear pairing over $\GF(2)$. Profile
linearity and evaluation linearity are its two
components. All six techniques are properties of this
pairing or operations on the spaces it connects:
\begin{itemize}[nosep]
  \item $1{+}1{=}0$: the field over which the pairing
    is defined.
  \item Leibniz: the boundary operator derived from the
    pairing's compatibility with the product.
  \item Linearity: the pairing is linear in the second
    argument.
  \item Strict: the pairing's product ($\wedge$) is
    strictly associative.
  \item SES: the pairing decomposes along the
    $\kappa$-split.
  \item Galois: the pairing is $GL$-invariant.
\end{itemize}

The six techniques are not six independent tools. They
are six facets of one bilinear pairing over $\GF(2)$,
and the 47 claims are 47 projections of this pairing
onto different parts of the mathematical landscape.


\subsection{Infinity is non-rotatability from inside}
\label{sec:infinity-nonrotatable}

The technique matrix upgrades the treatment of infinity
from axiomatic dissolution to structural necessity.

\emph{The five rotatable techniques are properties.}
$1{+}1{=}0$ is true on both sides of the $\kappa$-split.
Linearity holds inside and outside. Strict associativity,
the SES decomposition, the Galois group structure --- all
are symmetric across the inner/outer boundary. They can
be rotated through $S_3$ because they describe the
algebra, and the algebra is the same from every
perspective.

\emph{Leibniz is an operation, not a property.}
$\partial(R \wedge e_{n+1}) = R$ takes a cycle in
$\bigwedge(\kappa)$ and fills it with material from
$\bigwedge(V/\kappa)$. It is inherently directional:
from outer to inner, from complement to regime, from
unknown to known. It cannot be rotated because it IS
the asymmetry between perspectives. The other techniques
describe what is true on either side of the boundary.
Leibniz describes what happens when you cross it.

\emph{The experience of infinity from inside.}
From inside $\kappa$, each Leibniz crossing produces new
wedge products (the new $e_{n+1}$ wedges with everything
already in $\kappa$), which produce new residues, which
demand new crossings. The process is self-feeding. Each
filling creates the conditions for the next filling.
From inside, the supply of crossings looks inexhaustible.

From outside, each crossing reduces $\dim(V/\kappa)$ by
exactly 1. The Galois group shrinks. The process
terminates in $n$ steps.

The inner system cannot rotate Leibniz to see it as an
internal property, because Leibniz is the operation of
acquiring external resources. From inside, ``there is
always more to acquire'' is indistinguishable from ``the
supply is infinite,'' because the operation that would
resolve the distinction is the same operation that extends
the regime. Looking and extending are the same act.

\emph{Structural upgrade of Proposition
\ref{prop:indistinguish}.}
The original proof dissolves the actual/potential
distinction via the operationalist axiom: no discriminator
separates them, so they are identified. The technique
matrix provides a sharper statement:

\begin{quote}
The actual/potential distinction is unresolvable because
resolving it would require rotating the Leibniz technique
into a property of the regime. The Leibniz technique is
the unique element of the technique space that does not
admit rotation (\S\ref{sec:leibniz-exception}). The
non-rotatability is structural (proved from the technique
co-occurrence tree), not axiomatic (assumed by fiat).
\end{quote}

The connection: Leibniz is the filling operation. Filling
is $\kappa$-evolution. $\kappa$-evolution is what generates
the inner system's experience of inexhaustibility. You
cannot use Leibniz to determine whether the crossings
terminate, because using Leibniz IS crossing. The tool
and the question are the same act.

\emph{The halting tower.}
The reason the level-invariance holds at every level of
the halting tower (Remark \ref{rem:hyper-recursive}) is
not merely that each level faces the same structural
situation. It is that the operation connecting levels
(Leibniz filling at the meta-level) is the same
non-rotatable operation as at the object level. The tower
is self-similar because the technique that builds each
level is the one technique that cannot be recharacterized
as a property of the levels it builds.

\emph{The 3.1\% figure.}
At $n = 3$, only $8/256 = 3.1\%$ of Boolean functions are
admissible as primitive discriminators (linear functionals
over $\GF(2)$). This ratio $2^n / 2^{2^n}$ shrinks
exponentially: $0.024\%$ at $n = 4$, roughly $10^{-300}$
at $n = 10$. The primitive set becomes vanishingly small
while remaining fully expressive through composition.

This ratio is a \emph{figure of merit}, not a cost. The
smaller the primitive set, the more work the composition
mechanism ($\kappa$-evolution via Leibniz filling) does.
The more work the composition mechanism does, the more
algebraic invariants are maintained across every
operation. The more invariants maintained, the more
theorems follow structurally --- Russell exclusion,
foundation bounds, the bilinear form, the chain complex,
the Galois structure, all 47 claims.

A framework admitting 50\% of functions as primitives
would be expressively convenient and algebraically inert.
Nothing would be excluded structurally; every pathology
would need its own axiom. The theorems-per-axiom yield
would be low.

What the discrimination framework achieves is a
\emph{compression ratio for mathematical foundations}:
the primitive set is maximally small subject to the
constraint that composition recovers full expressiveness.
This is the same design principle as choosing
characteristic 2: pay in primitive expressiveness,
collect in structural yield. The 3.1\% and the $1+1=0$
are the same choice, made at different levels.
\subsection{Self-application}

The framework discriminates its own proofs by the same
structure it uses to discriminate data:
\begin{itemize}[nosep]
  \item Each claim has a technique profile in $\GF(2)^6$
    (a ``discriminator'' on the proof space).
  \item Claims with the same profile are
    ``$\kappa$-equivalent'' (proved by the same methods).
  \item The equivalence classes partition the paper by
    epistemological type: epistemic, constructive,
    coherent, geometric, operational, algebraic,
    universal, symmetric, higher, minimal.
  \item The rotation principle is the $S_3$ action on
    the proof space: the same symmetry that permutes
    $\{\kappa, \sigma, \tau\}$ permutes the three arms
    of the technique tree.
  \item The Galois group of the technique space acts on
    the ``unseen proofs'': alternate constructions are
    the Galois orbit of the original proof.
\end{itemize}

The $S_3$ acts on discriminators. It acts on data. It
acts on perspectives. And it acts on proofs about all
of the above. The technique matrix is the paper's own
exterior algebra, graded by proof complexity, with the
same Galois structure at the meta-level that the paper
describes at the object level.
